\chapter{The Soup That Was Free}

\section*{Told by}
Hag Morag (or so she called herself). She sat by a well, stirring a pot that had no fire under it. The soup was hot. The well was dry. The pot had been stirring itself, she said, for longer than the village had stood. She did not explain. She did not need to. A woman who sits by a dry well and stirs hot soup with no fire is not in the business of explanations.

\section*{The Tale}

A starving family came to a cottage. Not their own—they had lost their own to a fire, to a tax collector, to a winter that would not end. The story varied depending on who told it. What did not vary was the hunger. The mother's ribs showed through her dress. The father's hands shook when he knocked. The children did not cry. They had forgotten how.

A crone answered the door. Her eyes were the color of old milk. Her hands were the color of bark. She wore a shawl that had been mended so many times it was more thread than cloth.

\textit{``I have soup,''} she said. \textit{``It is free. The first bowl is always free. That is the rule. I did not make the rule. The rule was here before me. The rule will be here after.''}

They ate. The soup was rich, meaty, warm—thicker than any soup they had ever tasted, darker than any broth they had ever seen. It had chunks of something soft and something chewy and something that melted on the tongue like butter. They did not ask what was in it. Hungry people do not ask. Hungry people eat.

They ate until they were full. The children's cheeks flushed. The mother's hands stopped trembling. The father sat back in his chair and sighed—a long, slow breath that seemed to leave him lighter than he had been when he arrived.

The crone watched. Her eyes did not blink. \textit{``The soup is free,''} she said again. \textit{``The first bowl is free. But the seconds will cost you. The rule is clear. The bowl does not lie.''}

The family hesitated. The father said, \textit{``We have no coin. We have nothing. You can see our bones through our skin. What could we possibly have that you would want?''}

\textit{``I do not need coin,''} said the crone. \textit{``Coin is for merchants. I am not a merchant. I am a cook. I need a breath. One breath from each of you. Not your last breath—do not worry. I am not the Pale Shepherd. I do not take what you need to live. I take what you have to spare. Hold a breath for me. Just one. I will give it back when you next share a meal—a meal you cook yourselves, at your own hearth, with your own hands.''}

They agreed. What else could they do? They were full. They were grateful. They breathed into her kettle—one long, slow exhalation each, watching their breath fog the rim of the pot, watching the steam rise and curl and disappear into the crone's shawl. The kettle steamed. The soup bubbled. The crone nodded.

They left. That night, they dreamed they were drowning—not in water, but in something thicker, something that tasted of salt and old meat and the memory of hunger. They woke gasping, their chests tight, their breath shorter than it had been the day before.

They returned to the cottage. The crone was stirring the same pot. The fire was still absent. The soup was still hot.

\textit{``The debt is not due yet,''} she said. \textit{``The rule allows for patience. Eat again. The first bowl is still free. The second bowl will cost another breath. The third will cost two. The fourth will cost four. The rule is clear. The bowl does not lie.''}

They ate. They were not hungry. The soup had filled them, and the soup had not left. It sat in their bellies like a stone, warm and heavy and patient. They ate anyway. The father ate because he was afraid. The mother ate because she was tired. The children ate because the soup was warm and the night was cold and they had stopped asking questions.

The crone took another breath. The kettle grew heavier. The soup darkened.

For a year, they ate free soup once a week. Each time, the crone took a breath. Each time, the sleep grew deeper, the waking shorter, the dreams more vivid. They dreamed of drowning in soup. They dreamed of drowning in the kettle. They dreamed of the crone's eyes, old-milk white, watching them sink.

On the last night—the night of the forty-eighth bowl, the night when the father's breath came in rasps and the mother's steps dragged on the floor and the children slept for sixteen hours at a stretch—the father said, \textit{``We cannot pay. We have no breath left. The rule is clear. The bowl does not lie. What happens when we have nothing left to give?''}

The crone nodded. Her eyes, for the first time, softened. \textit{``Good,''} she said. \textit{``Now you are honest. The honesty is the payment I was waiting for. The soup was the toll. The breaths were the interest. The honesty is the principal. Go home. Cook your own meals. Breathe your own air. Dream your own dreams. The kettle will remember you. The kettle always remembers. But it will not follow. It does not need to. You know where it lives.''}

They went home. They felt lighter—not the lightness of a full belly, but the lightness of a released breath, of a held note finally sung, of a door that had been waiting to close. Their breath returned. Their sleep cleared. The dreams stopped.

They never ate free soup again. But sometimes, on winter nights, when the wind was right, they could smell it—rich and meaty and warm, drifting from the direction of the dry well. They did not follow the smell. They had learned. The soup was free. The seconds were not. The rule was clear. The bowl did not lie.

\section*{The Witch's Closing}
\textit{``Free soup is a hook. The hook is not the soup. The hook is the second bowl. The third bowl. The bowl you eat when you are not hungry. The bowl you eat because you are afraid of being hungry again. The crone does not need your coin. She needs your breath. And your breath, once given, is never quite your own again.''}

\section*{What Lurks in the Shadow of the Tale}

\subsection*{The Hungry Kettle}

A pot that produces unlimited nourishing soup—thick, dark, warm, satisfying in ways that have nothing to do with taste. The soup fills the belly. The soup warms the bones. The soup also takes. Each bowl consumed steals a small amount of breath, a sliver of life, a thread from the tapestry of the eater's waking hours. The soup is a loan. The interest compounds. The kettle does not lie.

\paragraph{Tags / Mechanics:}
\begin{itemize}
\item [HEAL] [CURSE] [SACRIFICE] [PRESSURE]
\item The soup restores 1 Supply per bowl. It also restores 1 Fatigue—but only temporarily. The eater feels refreshed for one hour. Then the Fatigue returns, along with an additional point of Fatigue that cannot be recovered by rest. The soup cheats. The soup always cheats.
\item Each bowl causes the eater to mark 1 Fatigue (which cannot be recovered by rest until the curse is lifted). This Fatigue stacks. There is no limit except the eater's willingness to keep eating.
\item \textbf{Clock: The Kettle's Patience [4]} — each bowl ticks once. When the timer fills, the eater must perform a favor for the hag (or a hag of her choosing—they know each other, they share notes, they compare debts). The favor will be strange. The favor will be difficult. The favor will not be dangerous, precisely—the hag does not need danger. She needs compliance. Refusal causes the accumulated Fatigue to become permanent. It cannot be healed by any means. It becomes part of the eater, like a scar, like a limp, like a name you cannot forget.
\item The kettle can be destroyed—shattered, buried, thrown into deep water. It will reform within a year. The hag is patient. The kettle is patient. The soup is always free.
\end{itemize}

\paragraph{Adventure Hook:}
A village has been living on free soup from a mysterious old woman who appears at the edge of the woods every night. The villagers are growing weak. They sleep too much. Their breaths come in rasps. Their children dream of drowning. The party must discover the price of the soup—and whether they can break the kettle without starving the village. The hag watches from the well. The hag does not interfere. The hag does not need to. The villagers will return on their own. They always return. The soup is free. The second bowl is not.

\section*{Colophon}
Well at a crossroads. The crone winked—a slow, deliberate motion, like a lizard closing one eye. \textit{``You did not eat,''} she said. \textit{``That is good. I would have taken your ink. Your ink is old. Your ink has seen things. Your ink would have made the soup bitter.''} She stirred the pot. The pot was empty. It had been empty the whole time. The soup had never been there. Or perhaps the soup had always been there, and the seeing was the eating, and the eating was the debt, and the debt was the story, and the story was the soup, and the soup was free. She did not explain. She did not need to.